%==================================================================
% LAMPIRAN
% Setiap lampiran dibuat dengan perintah \lampiran{Judul Lampiran}
% dan otomatis muncul di DAFTAR LAMPIRAN.
%==================================================================

\lampiran{Kode Program Pelatihan Jaringan Saraf Tiruan}

\noindent Berikut adalah kode program lengkap untuk melatih model
jaringan saraf tiruan yang digunakan pada penelitian ini.

\lstinputlisting[language=Python,
                 caption={Kode Program Pelatihan Model},
                 label={lst:pelatihan}]{kode/latih_model.py}

\clearpage

\lampiran{Lembar Observasi Kesegaran Ikan}

\noindent Formulir berikut digunakan untuk mencatat hasil pengamatan
organoleptik pada setiap objek penelitian.

\begin{spacing}{1}
\begin{tabularx}{\textwidth}{|>{\centering\arraybackslash}p{0.9cm}|X|>{\centering\arraybackslash}p{2.8cm}|>{\centering\arraybackslash}p{2.8cm}|}
\hline
\rowcolor{blue!12}
\textbf{No} & \textbf{Parameter Pengamatan} & \textbf{Skor (1--3)} & \textbf{Keterangan} \\
\hline
1 & Kejernihan mata ikan & & \\
\hline
2 & Warna insang & & \\
\hline
3 & Kekenyalan daging & & \\
\hline
4 & Bau & & \\
\hline
5 & Permukaan tubuh & & \\
\hline
\end{tabularx}
\end{spacing}

\vspace{2pt}
\noindent Keterangan skor: 1 = tidak segar, 2 = cukup segar, 3 = segar.
